\chapter{PHÂN TÍCH CẤU TRÚC}

\section{Giới thiệu giai đoạn}

\textbf{Phân tích cấu trúc} là giai đoạn thứ hai trong ba giai đoạn phân tích hướng đối tượng. Giai đoạn này trả lời: hệ thống được cấu thành từ những \textbf{lớp đối tượng} nào? Mỗi lớp có thuộc tính gì? Quan hệ giữa các lớp ra sao?

\begin{itemize}[leftmargin=1.5cm]
    \item \textbf{Biểu đồ chính}: Class Diagram (biểu đồ lớp).
    \item \textbf{Trước đó}: Ch.3 Phân tích chức năng -- Use Case (đã biết actor và mục tiêu).
    \item \textbf{Sau đó}: Ch.5 Phân tích hành vi -- Sequence (các đối tượng tương tác thế nào theo thời gian).
\end{itemize}

Từ đặc tả use case, các danh từ nghiệp vụ quan trọng được ứng viên thành lớp (ví dụ: hàng hóa, phiếu nhập, tồn kho, người dùng\ldots). Các lớp liên quan chặt được nhóm thành sơ đồ con để dễ đọc.

\section{Xác định các lớp chính}

\begin{longtable}{|p{3.5cm}|p{5cm}|p{4.5cm}|}
\hline
\textbf{Lớp} & \textbf{Vai trò} & \textbf{Use case liên quan} \\
\hline
\endfirsthead
\hline
\textbf{Lớp} & \textbf{Vai trò} & \textbf{Use case liên quan} \\
\hline
\endhead
NguoiDung & Tài khoản, vai trò, trạng thái & UC01--UC07 \\
\hline
HangHoa & Thông tin mặt hàng & UC13--UC17, phiếu \\
\hline
NhomHangHoa & Phân loại hàng & UC13 \\
\hline
DonViTinh & Đơn vị đo (cái, kg\ldots) & UC13 \\
\hline
NhaCungCap & Nhà cung cấp & UC08--UC12, UC23 \\
\hline
Kho & Kho hàng & UC18--UC22, phiếu, tồn \\
\hline
TonKho & Số lượng tồn theo kho--hàng & UC13, UC23--UC34 \\
\hline
PhieuNhapKho & Phiếu nhập & UC23--UC27 \\
\hline
ChiTietPhieuNhap & Dòng hàng trên phiếu nhập & UC23 \\
\hline
PhieuXuatKho & Phiếu xuất (ghi chú/lý do, không KH) & UC28--UC32 \\
\hline
ChiTietPhieuXuat & Dòng hàng trên phiếu xuất & UC28 \\
\hline
PhieuKiemKe & Phiếu kiểm kê & UC33 \\
\hline
ChiTietKiemKe & Dòng kiểm kê (thực tế, sổ, chênh) & UC33 \\
\hline
LichSuThaoTac & Log thao tác & Các UC ghi nhận \\
\hline
\end{longtable}

\section{Phân tích cấu trúc bằng biểu đồ lớp tổng thể}

Biểu đồ lớp tổng thể mô tả cấu trúc toàn hệ thống. Tách thành ba nhóm để tránh chồng chéo.

\subsection{Nhóm danh mục (Master Data)}

Gồm các lớp ít thay đổi: người dùng, hàng hóa, nhóm hàng, đơn vị tính, nhà cung cấp, kho, tồn kho.

\input{diagrams/tex/4.1_classdiagram_master}

Giải thích sơ đồ: Sơ đồ lớp danh mục mô tả ``xương sống'' dữ liệu tĩnh của hệ thống. \texttt{HangHoa} là trung tâm: thuộc một \texttt{NhomHangHoa}, gắn một \texttt{DonViTinh}, và có nhiều bản ghi \texttt{TonKho} (mỗi kho một dòng tồn). \texttt{Kho} quan hệ 1--n với \texttt{TonKho}. \texttt{NhaCungCap} dùng trên phiếu nhập; phiếu xuất dùng ghi chú người nhận (không có lớp khách hàng). \texttt{NguoiDung} đại diện tài khoản/vai trò. Các đường nối kèm số nhiều (1, n, 1..*) cho biết ràng buộc: ví dụ một hàng có nhiều tồn theo kho; nhiều hàng chung một nhóm.

\subsection{Nhóm nghiệp vụ nhập -- xuất}

Gồm phiếu nhập/xuất và chi tiết dòng hàng; quan hệ với NCC, KH, kho, hàng.

\input{diagrams/tex/4.2_classdiagram_nhapxuat}

Giải thích sơ đồ: Nhóm này mô tả cấu trúc \textbf{giao dịch}. \texttt{PhieuNhapKho} gắn \texttt{NhaCungCap}, \texttt{Kho}, \texttt{NguoiDung} (người lập) và có nhiều \texttt{ChiTietPhieuNhap}; mỗi chi tiết trỏ một \texttt{HangHoa}. \texttt{PhieuXuatKho} tương tự phiếu nhập nhưng không gắn NCC/KH -- có trường lý do/ghi chú. \texttt{TonKho} liên kết Kho--HangHoa để khi duyệt phiếu có thể tăng/giảm số lượng. Sơ đồ làm rõ: phiếu và chi tiết là cấu trúc master--detail; tồn là lớp bị cập nhật bởi hành vi duyệt (sẽ thấy rõ hơn ở sequence).

\subsection{Nhóm kiểm kê và lịch sử}

Gồm phiếu kiểm kê, chi tiết kiểm kê và log thao tác.

\input{diagrams/tex/4.3_classdiagram_kiemke_donhang}

Giải thích sơ đồ: \texttt{PhieuKiemKe} -- \texttt{ChiTietKiemKe} cũng theo mẫu master--detail; mỗi dòng lưu SL sổ, SL thực tế, chênh lệch. \texttt{LichSuThaoTac} ghi nhận ai làm gì (duyệt, thêm hàng\ldots), quan hệ với \texttt{NguoiDung}. Nhóm này bổ sung phần cấu trúc cho UC33 và audit, tách khỏi sơ đồ nhập xuất để tránh rối.

\section{Phân tích cấu trúc theo nhóm entity (biểu đồ lớp cắt lát)}

Thay vì vẽ lặp class cho từng UC CRUD (Thêm/Sửa/Xóa/Tìm/Xem DS cùng dùng một tập lớp), phần này \textbf{cắt lát theo nhóm entity}: mỗi nhóm chức năng một class diagram. Các UC cùng nhóm \textbf{chia sẻ cùng cấu trúc lớp}; khác biệt nằm ở \textbf{phương thức} gọi ra (thêm vs sửa vs xóa\ldots) và sẽ được thể hiện ở biểu đồ trình tự (Ch.5).

\begin{longtable}{|p{4.2cm}|p{3.8cm}|p{4.5cm}|}
\hline
\textbf{Nhóm entity} & \textbf{Use case} & \textbf{Sequence then chốt (Ch.5)} \\
\hline
\endfirsthead
\hline
\textbf{Nhóm entity} & \textbf{Use case} & \textbf{Sequence then chốt (Ch.5)} \\
\hline
\endhead
Tài khoản / xác thực & UC01--UC07 & UC01, UC03 \\
\hline
Nhà cung cấp & UC08--UC12 & UC08 \\
\hline
Hàng hóa & UC13--UC17 & UC13, UC16 \\
\hline
Kho & UC18--UC22 & UC18 \\
\hline
Phiếu nhập & UC23--UC27 & UC23 \\
\hline
Phiếu xuất & UC28--UC32 & UC28 \\
\hline
Kiểm kê & UC33 & UC33 \\
\hline
Báo cáo NXT & UC34 & UC34 \\
\hline
\end{longtable}

\subsection{Nhóm tài khoản / xác thực (UC01--UC07)}

\textbf{Mục tiêu nhóm}: đăng nhập/đăng xuất; QTV duy trì danh mục người dùng (thêm/sửa/xóa/tìm/xem DS).

\textbf{Lớp tham gia}: \texttt{NguoiDung}, \texttt{LichSuThaoTac}.

\textbf{Phương thức}: \texttt{dangNhap}/\texttt{dangXuat} (UC01/02); \texttt{themTaiKhoan} (UC03), \texttt{suaTaiKhoan} (UC04), \texttt{xoa}/\texttt{khoaMo} (UC05), \texttt{timKiem} (UC06), \texttt{layDanhSach} (UC07); \texttt{ghiLog}.

\begin{figure}[H]
\centering
\makebox[\textwidth][c]{%
\begin{tikzpicture}[scale=0.75, transform shape, every class/.style={text width=5.2cm}]

\begin{class}{NguoiDung}{-6,2}
\attribute{+ MaNguoiDung: String}
\attribute{+ TenDangNhap: String}
\attribute{+ MatKhau: String}
\attribute{+ HoTen: String}
\attribute{+ VaiTro: String}
\attribute{+ TrangThai: String}
\operation{+ dangNhap() / dangXuat()}
\operation{+ themTaiKhoan()}
\operation{+ suaTaiKhoan()}
\operation{+ xoa / khoaMoTaiKhoan()}
\operation{+ timKiemTaiKhoan()}
\operation{+ layDanhSachTaiKhoan()}
\end{class}

\begin{class}{LichSuThaoTac}{4,2}
\attribute{+ MaLog: String}
\attribute{+ MaNguoiDung: String}
\attribute{+ HanhDong: String}
\attribute{+ ThoiGian: DateTime}
\operation{+ ghiLog()}
\end{class}

\draw (NguoiDung) -- (LichSuThaoTac) node[midway, above, font=\tiny] {1..*};

\end{tikzpicture}
}
\caption{Biểu đồ lớp theo nhóm entity -- Tài khoản / xác thực (UC01--UC07)}
\label{fig:class-uc03}
\end{figure}


Giải thích: Cùng lớp \texttt{NguoiDung} phục vụ cả xác thực lẫn CRUD tài khoản. UC04--UC07 dùng lại cấu trúc này, chỉ khác operation. Sequence then chốt: UC01 (đăng nhập), UC03 (thêm TK).

\subsection{Nhóm nhà cung cấp (UC08--UC12)}

\textbf{Mục tiêu nhóm}: danh mục NCC phục vụ phiếu nhập -- đủ Thêm/Sửa/Xóa/Tìm/Xem DS.

\textbf{Lớp tham gia}: \texttt{NhaCungCap} (trung tâm), \texttt{NguoiDung} (actor), \texttt{PhieuNhapKho} (ràng buộc xóa), \texttt{LichSuThaoTac} (tùy chọn).

\textbf{Phương thức}: \texttt{themMoi} (UC08), \texttt{capNhat} (UC09), \texttt{xoa}/\texttt{ngungHoatDong} (UC10), \texttt{timKiem} (UC11), \texttt{layDanhSach} (UC12).

\begin{figure}[H]
\centering
\resizebox{\textwidth}{!}{%
\begin{tikzpicture}[
  every class/.style={text width=4.2cm, font=\small},
  node distance=1.2cm
]

% Hàng trên: actor | entity chính | ràng buộc
\begin{class}{NguoiDung}{-7.2,1.5}
\attribute{+ MaNguoiDung: String}
\attribute{+ VaiTro: String}
\end{class}

\begin{class}{NhaCungCap}{0,1.5}
\attribute{+ MaNhaCungCap: String}
\attribute{+ TenNhaCungCap: String}
\attribute{+ DiaChi: String}
\attribute{+ SoDienThoai: String}
\attribute{+ Email: String}
\attribute{+ TrangThai: String}
\operation{+ themMoi()}
\operation{+ capNhat()}
\operation{+ xoa() / ngungHoatDong()}
\operation{+ timKiem()}
\operation{+ layDanhSach()}
\end{class}

\begin{class}{PhieuNhapKho}{7.2,1.5}
\attribute{+ MaPhieuNhap: String}
\attribute{+ MaNhaCungCap: String}
\attribute{+ NgayNhap: DateTime}
\end{class}

% Hàng dưới: log (cách xa)
\begin{class}{LichSuThaoTac}{-7.2,-4.5}
\attribute{+ MaLog: String}
\attribute{+ HanhDong: String}
\attribute{+ ThoiGian: DateTime}
\end{class}

\draw (NguoiDung) -- (NhaCungCap) node[midway, above, font=\tiny] {thao tác};
\draw (NhaCungCap) -- (PhieuNhapKho) node[midway, above, font=\tiny] {1..*};
\draw (NguoiDung.south) -- (LichSuThaoTac.north) node[midway, left, font=\tiny] {1..*};

\end{tikzpicture}%
}
\caption{Biểu đồ lớp theo nhóm entity -- Nhà cung cấp (UC08--UC12)}
\label{fig:class-ncc}
\end{figure}


Giải thích: CRUD NCC thao tác trên cùng lớp \texttt{NhaCungCap}. Quan hệ tới \texttt{PhieuNhapKho} giải thích vì sao UC10 có thể chỉ \textit{ngừng} thay vì xóa cứng. Sequence then chốt: UC08; UC09--12 cùng luồng boundary--control--entity, đổi tên operation.

\subsection{Nhóm hàng hóa (UC13--UC17)}

\textbf{Mục tiêu nhóm}: danh mục hàng; khi thêm khởi tạo tồn; tìm/xem có thể kèm tồn.

\textbf{Lớp tham gia}: \texttt{HangHoa}, \texttt{NhomHangHoa}, \texttt{DonViTinh}, \texttt{TonKho}, \texttt{NguoiDung}.

\textbf{Phương thức}: \texttt{themHangHoa}+\texttt{khoiTaoTon} (UC13), \texttt{capNhatThongTin} (UC14), \texttt{xoa}/\texttt{ngung} (UC15), \texttt{timKiem} (UC16), \texttt{layDanhSach} (UC17).

\input{diagrams/tex/4.4_class_uc05}

Giải thích: \texttt{HangHoa} trung tâm; n..1 với nhóm và ĐVT; 1..* với \texttt{TonKho}. UC14--17 không đổi cấu trúc lớp. Sequence then chốt: UC13 (ghi + khởi tạo tồn), UC16 (chỉ đọc).

\subsection{Nhóm kho (UC18--UC22)}

\textbf{Mục tiêu nhóm}: danh mục kho (chỉ QL kho); tồn gắn theo từng kho.

\textbf{Lớp tham gia}: \texttt{Kho}, \texttt{TonKho}, \texttt{HangHoa}, \texttt{NguoiDung}.

\textbf{Phương thức}: \texttt{themMoi} (UC18), \texttt{capNhat} (UC19), \texttt{xoa}/\texttt{ngung} (UC20), \texttt{timKiem} (UC21), \texttt{layDanhSach} (UC22).

\begin{figure}[H]
\centering
\makebox[\textwidth][c]{%
\begin{tikzpicture}[scale=0.8, transform shape, every class/.style={text width=5.5cm}]

\begin{class}{Kho}{0,0}
\attribute{+ MaKho: String}
\attribute{+ TenKho: String}
\attribute{+ DiaChi: String}
\attribute{+ NguoiPhuTrach: String}
\attribute{+ TrangThai: String}
\operation{+ themMoi()}
\operation{+ capNhat()}
\operation{+ xoa() / ngungHoatDong()}
\operation{+ timKiem()}
\operation{+ layDanhSach()}
\end{class}

\begin{class}{TonKho}{6.5,0}
\attribute{+ MaKho: String}
\attribute{+ MaHangHoa: String}
\attribute{+ SoLuongTon: Integer}
\operation{+ layTon()}
\end{class}

\begin{class}{NguoiDung}{-6.5,0}
\attribute{+ MaNguoiDung: String}
\attribute{+ VaiTro: String}
\end{class}

\begin{class}{HangHoa}{6.5,-5}
\attribute{+ MaHangHoa: String}
\attribute{+ TenHangHoa: String}
\end{class}

\draw (Kho) -- (TonKho) node[midway, above, font=\tiny] {1..*};
\draw (TonKho) -- (HangHoa) node[midway, right, font=\tiny] {n..1};
\draw (NguoiDung) -- (Kho) node[midway, above, font=\tiny] {QL kho};

\end{tikzpicture}
}
\caption{Biểu đồ lớp theo nhóm entity -- Kho (UC18--UC22)}
\label{fig:class-kho}
\end{figure}


Giải thích: \texttt{Kho} 1..* \texttt{TonKho}; mỗi dòng tồn trỏ một \texttt{HangHoa}. Xóa kho bị chặn nếu còn tồn/phiếu. Sequence then chốt: UC18.

\subsection{Nhóm phiếu nhập (UC23--UC27)}

\textbf{Mục tiêu nhóm}: lập/sửa/xóa/tìm/xem phiếu nhập; duyệt tăng tồn.

\textbf{Lớp tham gia}: \texttt{PhieuNhapKho}, \texttt{ChiTietPhieuNhap}, \texttt{NhaCungCap}, \texttt{Kho}, \texttt{HangHoa}, \texttt{TonKho}, \texttt{NguoiDung}.

\textbf{Phương thức}: \texttt{lapPhieu}/\texttt{themChiTiet}/\texttt{duyet} (UC23); sửa/xóa khi chưa duyệt (UC24--25); tìm/xem DS (UC26--27); \texttt{TonKho.tangTon} khi duyệt.

\input{diagrams/tex/4.5_class_uc09}

Giải thích: Master--detail phiếu--chi tiết; duyệt mới đụng \texttt{TonKho}. UC24--27 cùng cấu trúc, khác bước trạng thái/truy vấn. Sequence then chốt: UC23.

\subsection{Nhóm phiếu xuất (UC28--UC32)}

\textbf{Mục tiêu nhóm}: lập/sửa/xóa/tìm/xem phiếu xuất; kiểm tra tồn; duyệt giảm tồn.

\textbf{Lớp tham gia}: \texttt{PhieuXuatKho}, \texttt{ChiTietPhieuXuat}, \texttt{Kho}, \texttt{HangHoa}, \texttt{TonKho}, \texttt{NguoiDung} (không NCC/KH).

\textbf{Phương thức}: \texttt{lapPhieu} + \texttt{kiemTraDu} (UC28); UC29--32 tương tự CRUD phiếu nhập; \texttt{giamTon} khi duyệt.

\begin{figure}[H]
\centering
\resizebox{0.98\textwidth}{!}{%
\begin{tikzpicture}[every class/.style={text width=3.3cm, font=\scriptsize}]

% Hàng trên
\begin{class}{Kho}{0,7}
\attribute{+ MaKho: String}
\attribute{+ TenKho: String}
\end{class}

\begin{class}{TonKho}{6.5,7}
\attribute{+ MaKho: String}
\attribute{+ MaHangHoa: String}
\attribute{+ SoLuongTon: Integer}
\operation{+ CapNhatTonKho()}
\operation{+ KiemTraTonKho()}
\end{class}

% Hàng dưới -- 4 lớp, khoảng cách đủ
\begin{class}{NguoiDung}{-9.5,0}
\attribute{+ MaNguoiDung: String}
\attribute{+ VaiTro: String}
\end{class}

\begin{class}{PhieuXuatKho}{-3.5,0}
\attribute{+ MaPhieuXuat: String}
\attribute{+ MaKho: String}
\attribute{+ LyDoXuat: String}
\attribute{+ GhiChu: String}
\attribute{+ NgayXuat: DateTime}
\attribute{+ TongTien: Decimal}
\operation{+ LapPhieuXuat()}
\operation{+ DuyetPhieuXuat()}
\end{class}

\begin{class}{ChiTietPhieuXuat}{3,0}
\attribute{+ MaPhieuXuat: String}
\attribute{+ MaHangHoa: String}
\attribute{+ SoLuong: Integer}
\attribute{+ DonGia: Decimal}
\end{class}

\begin{class}{HangHoa}{9,0}
\attribute{+ MaHangHoa: String}
\attribute{+ TenHangHoa: String}
\end{class}

\draw (NguoiDung) -- (PhieuXuatKho) node[midway, above, font=\tiny] {1..*};
\draw (PhieuXuatKho) -- (ChiTietPhieuXuat) node[midway, above, font=\tiny] {1..*};
\draw (ChiTietPhieuXuat) -- (HangHoa) node[midway, above, font=\tiny] {n..1};
\draw (PhieuXuatKho) -- (Kho) node[midway, left, font=\tiny] {n..1};
\draw (Kho) -- (TonKho) node[midway, above, font=\tiny] {1..*};
\draw (TonKho.south) -- (HangHoa.north) node[midway, right, font=\tiny] {n..1};

\end{tikzpicture}%
}
\caption{Biểu đồ lớp theo nhóm entity -- Phiếu xuất kho (UC28--UC32)}
\label{fig:class-uc10}
\end{figure}


Giải thích: Đối xứng nhóm nhập; \texttt{TonKho} vừa kiểm tra vừa giảm. Sequence then chốt: UC28.

\subsection{Nhóm kiểm kê (UC33)}

\textbf{Mục tiêu}: đối chiếu SL sổ / thực tế; duyệt thì gán lại tồn.

\textbf{Lớp tham gia}: \texttt{PhieuKiemKe}, \texttt{ChiTietKiemKe}, \texttt{Kho}, \texttt{HangHoa}, \texttt{TonKho}, \texttt{NguoiDung}.

\begin{figure}[H]
\centering
\resizebox{0.98\textwidth}{!}{%
\begin{tikzpicture}[every class/.style={text width=3.3cm, font=\scriptsize}]

% Hàng trên
\begin{class}{NguoiDung}{-9,7}
\attribute{+ MaNguoiDung: String}
\attribute{+ VaiTro: String}
\end{class}

\begin{class}{Kho}{-3,7}
\attribute{+ MaKho: String}
\attribute{+ TenKho: String}
\end{class}

% Hàng giữa
\begin{class}{PhieuKiemKe}{-3,0}
\attribute{+ MaPhieuKK: String}
\attribute{+ MaKho: String}
\attribute{+ NgayKiemKe: DateTime}
\attribute{+ TrangThai: String}
\operation{+ taoPhieu()}
\operation{+ nhapSoThucTe()}
\operation{+ guiDuyet()}
\operation{+ duyetDieuChinh()}
\end{class}

\begin{class}{ChiTietKiemKe}{3.5,0}
\attribute{+ MaPhieuKK: String}
\attribute{+ MaHangHoa: String}
\attribute{+ SLSoSach: Integer}
\attribute{+ SLThucTe: Integer}
\attribute{+ ChenhLech: Integer}
\end{class}

\begin{class}{HangHoa}{9.5,0}
\attribute{+ MaHangHoa: String}
\attribute{+ TenHangHoa: String}
\end{class}

% Hàng dưới
\begin{class}{TonKho}{3.5,-6}
\attribute{+ MaKho: String}
\attribute{+ MaHangHoa: String}
\attribute{+ SoLuongTon: Integer}
\operation{+ ganTheoKiemKe()}
\end{class}

\draw (NguoiDung) -- (PhieuKiemKe) node[midway, left, font=\tiny] {1..*};
\draw (Kho) -- (PhieuKiemKe) node[midway, right, font=\tiny] {n..1};
\draw (PhieuKiemKe) -- (ChiTietKiemKe) node[midway, above, font=\tiny] {1..*};
\draw (ChiTietKiemKe) -- (HangHoa) node[midway, above, font=\tiny] {n..1};
\draw (ChiTietKiemKe) -- (TonKho) node[midway, right, font=\tiny] {đối chiếu};

\end{tikzpicture}%
}
\caption{Biểu đồ lớp theo nhóm entity -- Kiểm kê (UC33)}
\label{fig:class-kiemke}
\end{figure}


Giải thích: Master--detail kiểm kê; \texttt{ganTheoKiemKe} chỉ sau duyệt. Sequence: UC33.

\subsection{Nhóm báo cáo NXT (UC34)}

\textbf{Mục tiêu}: tổng hợp nhập--xuất--tồn theo kỳ/kho/hàng (chỉ đọc).

\textbf{Lớp tham gia}: lớp báo cáo (logic tổng hợp) + đọc \texttt{PhieuNhapKho}, \texttt{PhieuXuatKho}, \texttt{TonKho}, \texttt{HangHoa}, \texttt{Kho}.

\begin{figure}[H]
\centering
\makebox[\textwidth][c]{%
\begin{tikzpicture}[scale=0.8, transform shape, every class/.style={text width=5.2cm}]

\begin{class}{BaoCaoNXT}{0,2}
\attribute{+ TuNgay: Date}
\attribute{+ DenNgay: Date}
\attribute{+ MaKho: String}
\attribute{+ MaHang: String}
\operation{+ tongHopNhap()}
\operation{+ tongHopXuat()}
\operation{+ layTonHienTai()}
\operation{+ xuatExcel()}
\end{class}

\begin{class}{PhieuNhapKho}{-6.5,0}
\attribute{+ MaPhieuNhap: String}
\attribute{+ NgayNhap: DateTime}
\attribute{+ TrangThai: String}
\end{class}

\begin{class}{PhieuXuatKho}{6.5,0}
\attribute{+ MaPhieuXuat: String}
\attribute{+ NgayXuat: DateTime}
\attribute{+ TrangThai: String}
\end{class}

\begin{class}{TonKho}{-3.5,-4.5}
\attribute{+ MaKho: String}
\attribute{+ MaHangHoa: String}
\attribute{+ SoLuongTon: Integer}
\end{class}

\begin{class}{HangHoa}{3.5,-4.5}
\attribute{+ MaHangHoa: String}
\attribute{+ TenHangHoa: String}
\end{class}

\begin{class}{Kho}{0,-4.5}
\attribute{+ MaKho: String}
\attribute{+ TenKho: String}
\end{class}

\draw (BaoCaoNXT) -- (PhieuNhapKho) node[midway, left, font=\tiny] {đọc};
\draw (BaoCaoNXT) -- (PhieuXuatKho) node[midway, right, font=\tiny] {đọc};
\draw (BaoCaoNXT) -- (TonKho) node[midway, left, font=\tiny] {đọc};
\draw (TonKho) -- (HangHoa) node[midway, above, font=\tiny] {};
\draw (TonKho) -- (Kho) node[midway, above, font=\tiny] {};

\end{tikzpicture}
}
\caption{Biểu đồ lớp theo nhóm entity -- Báo cáo nhập--xuất--tồn (UC34)}
\label{fig:class-baocao}
\end{figure}


Giải thích: Không tạo entity giao dịch mới; chỉ \textbf{đọc} phiếu đã duyệt và tồn. Sequence: UC34.

\section{Bảng mô tả thuộc tính các lớp}

\begin{longtable}{|p{3cm}|p{5.5cm}|p{4.5cm}|}
\hline
\textbf{Lớp} & \textbf{Thuộc tính chính} & \textbf{Ghi chú} \\
\hline
\endfirsthead
\hline
\textbf{Lớp} & \textbf{Thuộc tính chính} & \textbf{Ghi chú} \\
\hline
\endhead
NguoiDung & MaNguoiDung, TenDangNhap, MatKhau, HoTen, VaiTro, TrangThai & Phân quyền \\
\hline
HangHoa & MaHangHoa, TenHangHoa, MaNhom, MaDVT, GiaNhap, GiaXuat, TrangThai & Trung tâm danh mục \\
\hline
NhomHangHoa & MaNhomHang, TenNhomHang, MoTa & Phân loại \\
\hline
DonViTinh & MaDonViTinh, TenDonViTinh & Cái, kg, thùng\ldots \\
\hline
NhaCungCap & MaNCC, Ten, DiaChi, SDT, Email, NguoiLienHe & Phục vụ nhập \\
\hline
Kho & MaKho, TenKho, DiaChi, NguoiPhuTrach, TrangThai & Nhiều kho \\
\hline
TonKho & MaKho, MaHangHoa, SoLuongTon, NgayCapNhat & Khóa phức hợp \\
\hline
PhieuNhapKho & MaPhieu, MaNCC, MaKho, MaNV, NgayNhap, TongTien, TrangThai & Chờ duyệt/Đã duyệt \\
\hline
ChiTietPhieuNhap & MaPhieu, MaHang, SoLuong, DonGia, ThanhTien & Dòng nhập \\
\hline
PhieuXuatKho & MaPhieu, MaKho, MaNV, NgayXuat, LyDoXuat, GhiChu, TongTien, TrangThai & Không MaKH \\
\hline
ChiTietPhieuXuat & MaPhieu, MaHang, SoLuong, DonGia, ThanhTien & Dòng xuất \\
\hline
PhieuKiemKe & MaPhieu, MaKho, MaNV, NgayKiem, TrangThai & Kiểm kê \\
\hline
ChiTietKiemKe & MaPhieu, MaHang, SLThucTe, SLSoSach, ChenhLech & Điều chỉnh tồn \\
\hline
LichSuThaoTac & MaLog, MaNguoiDung, HanhDong, ThoiGian, DoiTuong & Audit \\
\hline
\end{longtable}

\section{Giải thích các phương thức của lớp}

Ngoài thuộc tính (dữ liệu), mỗi lớp có \textbf{phương thức} (hành vi/khả năng xử lý). Dưới đây mô tả các phương thức chính gắn với use case đã phân tích.

\subsection{Lớp NguoiDung}

\begin{longtable}{|p{4.2cm}|p{8.5cm}|}
\hline
\textbf{Phương thức} & \textbf{Mô tả / khi nào dùng} \\
\hline
\endfirsthead
\hline
\textbf{Phương thức} & \textbf{Mô tả / khi nào dùng} \\
\hline
\endhead
\texttt{dangNhap(ten, matKhau)} & Kiểm tra tên đăng nhập và mật khẩu; trả về thành công/thất bại. Dùng cho UC01. \\
\hline
\texttt{dangXuat()} & Hủy phiên làm việc hiện tại. Dùng cho UC02. \\
\hline
\texttt{themTaiKhoan(thongTin)} & Tạo tài khoản mới (mã hóa MK, gán vai trò). Dùng cho UC03. \\
\hline
\texttt{suaTaiKhoan(ma, thongTin)} & Cập nhật họ tên, email, vai trò, (mật khẩu nếu đổi). Dùng cho UC04. \\
\hline
\texttt{xoaTaiKhoan(ma)} / \texttt{khoaMoTaiKhoan(ma, trangThai)} & Xóa nếu chưa có phiếu/log; nếu đã có thì chỉ khóa (UC05). \\
\hline
\texttt{timKiemTaiKhoan(tuKhoa, boLoc)} & Lọc theo tên đăng nhập / họ tên / vai trò / trạng thái. UC06. \\
\hline
\texttt{layDanhSachTaiKhoan()} & Tải danh sách (phân trang) -- UC07. \\
\hline
\texttt{kiemTraQuyen(hanhDong)} & Kiểm tra vai trò có được phép thực hiện hành động (lập phiếu, duyệt\ldots). \\
\hline
\end{longtable}

\subsection{Lớp HangHoa}

\begin{longtable}{|p{4.2cm}|p{8.5cm}|}
\hline
\textbf{Phương thức} & \textbf{Mô tả / khi nào dùng} \\
\hline
\endfirsthead
\hline
\textbf{Phương thức} & \textbf{Mô tả / khi nào dùng} \\
\hline
\endhead
\texttt{themHangHoa(thongTin)} & Lưu hàng mới; gọi khởi tạo tồn kho. UC13. \\
\hline
\texttt{capNhatThongTin(ma, thongTin)} & Sửa tên, giá, nhóm, mô tả. UC14. \\
\hline
\texttt{ngungSuDung(ma)} / \texttt{xoa(ma)} & Xóa hoặc ngừng (nếu đã có phiếu). UC15. \\
\hline
\texttt{timKiem(tuKhoa, boLoc)} & Tìm theo mã/tên/nhóm. UC16. \\
\hline
\texttt{layDanhSach()} & Xem DS / phân trang. UC17. \\
\hline
\texttt{layThongTin(ma)} & Lấy chi tiết một mặt hàng. \\
\hline
\end{longtable}

\subsection{Lớp TonKho}

\begin{longtable}{|p{4.5cm}|p{8.2cm}|}
\hline
\textbf{Phương thức} & \textbf{Mô tả / khi nào dùng} \\
\hline
\endfirsthead
\hline
\textbf{Phương thức} & \textbf{Mô tả / khi nào dùng} \\
\hline
\endhead
\texttt{khoiTaoTon(maKho, maHang)} & Tạo dòng tồn = 0 khi thêm hàng mới (UC13). \\
\hline
\texttt{tangTon(maKho, maHang, sl)} & Cộng số lượng khi \textbf{duyệt phiếu nhập} (UC23). \\
\hline
\texttt{giamTon(maKho, maHang, sl)} & Trừ số lượng khi \textbf{duyệt phiếu xuất} (UC28); chỉ gọi nếu đã kiểm tra đủ. \\
\hline
\texttt{kiemTraDu(maKho, maHang, sl)} & Trả về true/false -- tồn có đủ để xuất không. UC28. \\
\hline
\texttt{ganTheoKiemKe(maKho, maHang, slThuc)} & Gán tồn = số thực tế khi duyệt kiểm kê (UC33). \\
\hline
\texttt{layTon(maKho, maHang)} & Đọc số lượng hiện tại (tìm kiếm, báo cáo UC16, UC34). \\
\hline
\end{longtable}

\subsection{Lớp PhieuNhapKho}

\begin{longtable}{|p{4.2cm}|p{8.5cm}|}
\hline
\textbf{Phương thức} & \textbf{Mô tả / khi nào dùng} \\
\hline
\endfirsthead
\hline
\textbf{Phương thức} & \textbf{Mô tả / khi nào dùng} \\
\hline
\endhead
\texttt{lapPhieu(thongTinChung)} & Tạo phiếu mới (NCC, kho, ngày, người lập); trạng thái ban đầu. UC23. \\
\hline
\texttt{themChiTiet(dongHang)} & Thêm dòng hàng (mã hàng, SL, đơn giá); tính thành tiền. \\
\hline
\texttt{tinhTongTien()} & Cộng thành tiền các dòng $\rightarrow$ \texttt{TongTien}. \\
\hline
\texttt{guiDuyet()} & Đổi trạng thái sang ``Chờ duyệt''. \\
\hline
\texttt{duyetPhieu()} & QL duyệt: ``Đã duyệt''; gọi \texttt{TonKho.tangTon} từng dòng; ghi log. \\
\hline
\texttt{tuChoiPhieu(lyDo)} & ``Từ chối''; \textbf{không} đổi tồn; lưu lý do. \\
\hline
\end{longtable}

\subsection{Lớp PhieuXuatKho}

\begin{longtable}{|p{4.2cm}|p{8.5cm}|}
\hline
\textbf{Phương thức} & \textbf{Mô tả / khi nào dùng} \\
\hline
\endfirsthead
\hline
\textbf{Phương thức} & \textbf{Mô tả / khi nào dùng} \\
\hline
\endhead
\texttt{lapPhieu(thongTinChung)} & Tạo phiếu xuất (KH, kho, ngày). UC28. \\
\hline
\texttt{themChiTiet(dongHang)} & Thêm dòng; trước đó gọi \texttt{TonKho.kiemTraDu}; đủ mới cho thêm. \\
\hline
\texttt{tinhTongTien()} & Tính tổng giá trị phiếu xuất. \\
\hline
\texttt{guiDuyet()} & Chuyển ``Chờ duyệt''. \\
\hline
\texttt{duyetPhieu()} & ``Đã duyệt''; gọi \texttt{TonKho.giamTon} từng dòng; ghi log. \\
\hline
\texttt{tuChoiPhieu(lyDo)} & Từ chối; tồn không đổi. \\
\hline
\end{longtable}

\subsection{Lớp PhieuKiemKe}

\begin{longtable}{|p{4.5cm}|p{8.2cm}|}
\hline
\textbf{Phương thức} & \textbf{Mô tả / khi nào dùng} \\
\hline
\endfirsthead
\hline
\textbf{Phương thức} & \textbf{Mô tả / khi nào dùng} \\
\hline
\endhead
\texttt{taoPhieu(maKho, ngay)} & Tạo phiếu kiểm kê; nạp danh sách hàng + SL sổ từ \texttt{TonKho}. UC33. \\
\hline
\texttt{nhapSoThucTe(maHang, sl)} & Ghi số thực tế; tính chênh lệch = thực tế $-$ sổ. \\
\hline
\texttt{guiDuyet()} & Gửi QL xem xét. \\
\hline
\texttt{duyetDieuChinh()} & QL duyệt: gọi \texttt{TonKho.ganTheoKiemKe} cho từng dòng lệch. \\
\hline
\end{longtable}

\subsection{Lớp NhaCungCap / Kho (và các danh mục tương tự)}

\begin{longtable}{|p{4.2cm}|p{8.5cm}|}
\hline
\textbf{Phương thức} & \textbf{Mô tả / khi nào dùng} \\
\hline
\endfirsthead
\hline
\textbf{Phương thức} & \textbf{Mô tả / khi nào dùng} \\
\hline
\endhead
\texttt{themMoi(thongTin)} & Thêm bản ghi mới: NCC (UC08) hoặc Kho (UC18). \\
\hline
\texttt{capNhat(ma, thongTin)} & Sửa thông tin danh mục -- NCC dùng cho UC09; hàng UC14; kho UC19. \\
\hline
\texttt{xoa(ma)} / \texttt{ngungHoatDong(ma)} & Xóa cứng nếu chưa có phiếu; nếu đã có phiếu nhập thì chỉ ngừng hoạt động (UC10); hàng UC15; kho UC20. \\
\hline
\texttt{timKiem(tuKhoa)} & Tra cứu theo mã/tên/SĐT -- NCC UC11; hàng UC16; kho UC21. \\
\hline
\texttt{layDanhSach()} & Tải danh sách (có phân trang) -- NCC UC12; hàng UC17; kho UC22; phiếu UC27/UC32. \\
\hline
\end{longtable}

\subsection{Lớp LichSuThaoTac}

\begin{longtable}{|p{4.5cm}|p{8.2cm}|}
\hline
\textbf{Phương thức} & \textbf{Mô tả / khi nào dùng} \\
\hline
\endfirsthead
\hline
\textbf{Phương thức} & \textbf{Mô tả / khi nào dùng} \\
\hline
\endhead
\texttt{ghiLog(nguoiDung, hanhDong, doiTuong, moTa)} & Ghi nhận ai làm gì, lúc nào (đăng nhập, duyệt phiếu, thêm hàng\ldots). \\
\hline
\end{longtable}

\subsection{Liên hệ phương thức với use case (tóm tắt)}

\begin{itemize}[leftmargin=1.5cm]
    \item \textbf{UC01/UC02}: \texttt{NguoiDung.dangNhap} / \texttt{dangXuat}; có thể kèm \texttt{LichSuThaoTac.ghiLog}.
    \item \textbf{UC03}: \texttt{themTaiKhoan}; \textbf{UC04}: \texttt{suaTaiKhoan}; \textbf{UC05}: \texttt{xoaTaiKhoan}/\texttt{khoaMoTaiKhoan}; \textbf{UC06}: \texttt{timKiemTaiKhoan}; \textbf{UC07}: \texttt{layDanhSachTaiKhoan}.
    \item \textbf{UC08--UC12}: \texttt{NhaCungCap} them/sua/xoa/tim/xem DS.
    \item \textbf{UC13--UC17}: \texttt{HangHoa} them (+\texttt{khoiTaoTon})/sua/xoa/tim/xem DS.
    \item \textbf{UC18--UC22}: \texttt{Kho} them/sua/xoa/tim/xem DS.
    \item \textbf{UC23--UC27}: \texttt{PhieuNhapKho} lập/sửa/xóa/tìm/xem (+ duyệt, \texttt{tangTon}).
    \item \textbf{UC28--UC32}: \texttt{PhieuXuatKho} tương tự (+ \texttt{kiemTraDu}, \texttt{giamTon}).
    \item \textbf{UC33}: \texttt{PhieuKiemKe.*} + \texttt{ganTheoKiemKe}.
    \item \textbf{UC34}: đọc/tổng hợp tồn và phiếu (báo cáo).
\end{itemize}

\section{Tóm tắt quan hệ quan trọng}

\begin{itemize}[leftmargin=1.5cm]
    \item \textbf{1 -- n}: NguoiDung--Phieu; NhaCungCap--PhieuNhap; Phieu--ChiTiet; Kho--TonKho; HangHoa--TonKho.
    \item \textbf{n -- 1}: ChiTiet--HangHoa; HangHoa--NhomHangHoa; HangHoa--DonViTinh.
    \item \textbf{Cập nhật tồn}: chỉ qua phương thức \texttt{tangTon} / \texttt{giamTon} / \texttt{ganTheoKiemKe} khi duyệt phiếu tương ứng.
\end{itemize}

\section{Kết luận chương}

Chương 4 hoàn thành \textbf{phân tích cấu trúc} bằng biểu đồ lớp: lớp tổng thể (danh mục, nhập--xuất, kiểm kê) và \textbf{8 cắt lát theo nhóm entity} (tài khoản, NCC, hàng, kho, phiếu nhập, phiếu xuất, kiểm kê, báo cáo) phủ toàn bộ UC01--UC34. Thuộc tính, phương thức và quan hệ đã nêu; phương thức then chốt xuất hiện lại trên \textbf{biểu đồ trình tự} (Ch.5).
